\section{Cómo leer una scaling law: fit, residual y compute allocation}

La sección anterior explicó la ruta de investigación; esta pasa a las fórmulas y a las decisiones. Las scaling laws habituales de los modelos de lenguaje expresan la cross-entropy loss como un decaimiento en ley de potencias respecto de los recursos más un término irreducible. Lo importante no es memorizar el exponent, sino entender cómo cada supuesto del fit, el rango de los datos y el error afectan una compute allocation que cuesta decenas de millones o más.

\subsection{Power law y pérdida irreducible}

La \term{cross-entropy loss} (pérdida de entropía cruzada) mide la probabilidad logarítmica negativa que el modelo asigna al siguiente token real. Una forma simplificada de compute scaling es
\[
L(C)=L_{\infty}+aC^{-\alpha},
\]
donde $C$ representa el compute de entrenamiento, $L(C)$ es la loss correspondiente, $L_{\infty}$ es el término irreducible bajo los supuestos actuales de datos y tarea, $a$ es el coeficiente de escala y $\alpha>0$ es la velocidad con la que decaen los rendimientos. En una gráfica log-log la curva puede parecer casi una recta, pero $L_{\infty}$, el ruido y los recipe changes afectan de manera considerable la extrapolación.

\lecturefigure{03-scaling-law-fit.png}{Una scaling law es un fit con residual, confidence y límites de validez, no una recta eterna.}{Kaplan et al. 2020; subtítulos locales 06:20--11:50; redibujo conceptual.}

\begin{knowledgebox}{Lectura de la figura: primero los residuos, luego la pendiente}
Los axes deben indicar si se trata de FLOPs, tokens, parameters o wall-clock; el fit debe reportar el intervalo de entrenamiento y la confidence; el residual muestra si la architecture, los data o la optimization se desvían de forma sistemática; solo entonces la decision puede ajustar el budget, la model/data mix o detener la extrapolación. Si las runs más grandes caen una tras otra del mismo lado de la línea de predicción, la información más importante quizá no sea «el modelo va bastante bien», sino que la law actual ya no se ajusta.
\end{knowledgebox}

\begin{warningbox}{La mejora en benchmarks y el loss scaling no son la misma función}
Un descenso suave de la average loss puede corresponder a efectos de umbral repentinos en algunas downstream tasks, o puede no mejorar la tarea objetivo. La evaluación de capacidades debe abarcar entornos reales, herramientas y elicitation; una sola curva de loss de preentrenamiento no sustituye la calidad de post-training, la seguridad ni la serving economics.
\end{warningbox}

\subsection{El compute multiplier es una combinación, no una constante secreta}

En clase se agrupan las mejoras de architecture, data, optimizer y systems bajo el nombre de compute multiplier: obtener más effective capability con los mismos nominal FLOPs. Mantenerlas en secreto tiene razones competitivas, pero internamente deben poder reproducirse. Cada multiplier debe contar con evidencia de baseline, ablation, interaction y transfer, para evitar que en el portfolio varios cambios se anulen entre sí y aun así todos se declaren exitosos.

\lecturefigure{04-compute-allocation.png}{La compute allocation invierte a la vez en cuatro tipos de multiplier: model, data, optimization y systems.}{Subtítulos locales 12:00--15:20; redibujo conceptual.}

\begin{knowledgebox}{Lectura de la figura: las ganancias se atribuyen según el effective compute}
La architecture puede mejorar la capacidad expresiva o la eficiencia por dispersión; los data mejoran la signal-to-noise; la optimization mejora la estabilidad y la convergencia; los systems mejoran la hardware utilization. Al final hay que comparar la quality con wall-clock fijo, dólares fijos, energy fija o FLOPs fijos, en lugar de reportar solo un microbenchmark. Además, un multiplier puede funcionar solo a cierta scale, por lo que es necesario verificar su transfer.
\end{knowledgebox}

\teachervoice{El ponente describe un frontier lab como una colaboración estrecha entre researchers e engineers, no como un lugar donde los investigadores proponen ideas y los ingenieros hacen el «trabajo menor». La scaling law les da a ambos un lenguaje experimental común: la investigación modifica la recipe y la ingeniería garantiza que cada run sea comparable, recuperable y medible.}

\subsection{Resumen de la sección}

La scaling law es una herramienta para decidir sobre recursos. Lo correcto es combinar el fit, el residual, la uncertainty y la ablation de los multipliers; lo incorrecto es tratar un exponente histórico como una constante natural o equiparar directamente una mejora en el throughput del sistema con una mejora en la capacidad del modelo.
